\documentclass[11pt]{article}

\usepackage[a4paper,margin=24mm]{geometry}
\usepackage[T1]{fontenc}
\usepackage[utf8]{inputenc}
\usepackage{lmodern}
\usepackage{microtype}
\usepackage{amsmath,amssymb,mathtools,bm}
\usepackage{booktabs,longtable,array}
\usepackage{xcolor}
\usepackage{listings}
\usepackage{enumitem}
\usepackage{float}
\usepackage{hyperref}

\hypersetup{
  colorlinks=true,
  linkcolor=blue!55!black,
  urlcolor=blue!55!black,
  citecolor=blue!55!black,
  pdftitle={Detailed Implementation Notes for VULETTA and VUMPS in pyqed}
}

\lstdefinestyle{pyqedpython}{
  language=Python,
  basicstyle=\ttfamily\small,
  keywordstyle=\color{blue!60!black},
  commentstyle=\color{green!40!black},
  stringstyle=\color{red!55!black},
  showstringspaces=false,
  breaklines=true,
  columns=fullflexible,
  frame=single,
  framerule=0.3pt
}

\newcommand{\Id}{\mathbb{I}}
\newcommand{\CC}{\mathbb{C}}
\newcommand{\R}{\mathbb{R}}
\newcommand{\E}{\mathcal{E}}
\newcommand{\Htwo}{\mathcal{H}}
\newcommand{\Res}{\mathcal{R}}
\newcommand{\Proj}{\mathcal{P}}
\newcommand{\Qproj}{\mathcal{Q}}
\newcommand{\Tr}{\operatorname{Tr}}
\newcommand{\rank}{\operatorname{rank}}
\newcommand{\diag}{\operatorname{diag}}
\newcommand{\vecop}{\operatorname{vec}}
\newcommand{\reshape}{\operatorname{reshape}}
\newcommand{\normalize}{\operatorname{normalize}}
\newcommand{\RePart}{\operatorname{Re}}
\newcommand{\argmin}{\operatorname*{arg\,min}}
\newcommand{\norm}[1]{\left\lVert #1\right\rVert}
\newcommand{\inner}[2]{\left\langle #1,#2\right\rangle}
\newcommand{\code}[1]{\texttt{\detokenize{#1}}}

\title{Detailed Implementation Notes for the Current \texttt{pyqed} VULETTA Solver\\
\large Exact LETTA--MPS Embeddings, Conditional Gauge Fixing, Analytic Gradients,
Tangent Updates, and Comparison with VUMPS}
\author{Code-grounded technical documentation}
\date{Source state: 27 August 2026}

\begin{document}
\maketitle

\begin{abstract}
This note documents the implementation in \path{pyqed/_vuletta} and compares it
line by line with the one-site VUMPS implementation in \path{pyqed/_vumps}.
The current VULETTA algorithm is a direct variational optimizer for a uniform
nearest-neighbor leg-tied tensor.  It uses an exact sparse MPS embedding only
to contract the infinite state.  It does not optimize an unrestricted MPS and
does not project an MPS back into the LETTA manifold.  The active solver first
constructs a physical-sector-conditioned mixed-canonical form, removes the
sector-dependent virtual gauge from the tangent space, whitens the remaining
tangent coordinates with the supported right metric, evaluates an analytic
transfer-response gradient, and performs an Armijo step whose retraction must
remain canonical and injective.  By contrast, the VUMPS implementation
optimizes an unrestricted one-site MPS through center-tensor and center-matrix
eigenproblems followed by polar gauge matching.  Every formula below is tied
to the current source code; this is an implementation description, not a
claim that the LETTA paper itself derived this optimizer.
\end{abstract}

\tableofcontents
\clearpage

\section{Purpose, scope, and the most important distinction}

Consider an infinite one-dimensional Hamiltonian with a translation-invariant
nearest-neighbor density,
\begin{equation}
  H=\sum_{n\in\mathbb{Z}} h_{n,n+1},
  \qquad
  h_{ij,kl}=\langle i,j|h|k,l\rangle,
  \label{eq:hamiltonian}
\end{equation}
where every local physical index takes values in
$\{0,\ldots,d-1\}$.  Both solvers minimize the thermodynamic energy density
\begin{equation}
  e=\lim_{N\rightarrow\infty}
  \frac{\langle\Psi_N|H_N|\Psi_N\rangle}
       {N\langle\Psi_N|\Psi_N\rangle}.
  \label{eq:energy-density-definition}
\end{equation}
They differ in the variational manifold and in the stationarity equation.

The current VULETTA implementation optimizes only
\begin{equation}
  T\in\CC^{D\times d\times d\times D},
  \qquad T^{p,q}_{\alpha\beta}
  =\code{tensor[alpha,p,q,beta]}.
  \label{eq:letta-shape}
\end{equation}
The two physical legs $p$ and $q$ are tied to neighboring sites.  The exact
MPS tensor introduced below is a contraction identity for this $T$; its zero
pattern is never released.  Therefore the algorithm is not VUMPS with a
larger bond dimension, and it is not an optimize-then-project scheme.

The VUMPS implementation instead optimizes an ordinary uniform MPS tensor
\begin{equation}
  A\in\CC^{D\times d\times D},
  \qquad A^s_{\alpha\beta}=\code{A[alpha,s,beta]},
  \label{eq:mps-shape}
\end{equation}
using mixed-canonical center equations.  Sections~\ref{sec:vuletta}--
\ref{sec:vuletta-solver} derive VULETTA; Section~\ref{sec:vumps} derives VUMPS;
Section~\ref{sec:comparison} gives a direct comparison.

\section{Notation ledger and array conventions}
\label{sec:notation}

\begingroup
\small
\begin{longtable}{p{0.19\textwidth}p{0.26\textwidth}p{0.47\textwidth}}
\toprule
Symbol & Code object or shape & Meaning \\
\midrule
\endhead
$d$ & physical dimension & Number of states on one site. \\
$D$ & LETTA bond dimension & Dimension of $\alpha,\beta$ in $T^{p,q}_{\alpha\beta}$. \\
$\chi=dD$ & \code{effective_bond_dim} & Bond dimension of the exact structured MPS used for LETTA contractions. \\
$p,q$ & $0,\ldots,d-1$ & Previous and current physical state in a LETTA block. \\
$s,t,u,v$ & $0,\ldots,d-1$ & Generic MPS or Hamiltonian physical indices. \\
$\alpha,\beta$ & $0,\ldots,D-1$ & LETTA virtual indices. \\
$a,b,c,e$ & $0,\ldots,\chi-1$ & Composite structured-MPS virtual indices. \\
$T^{p,q}$ & $D\times D$ & LETTA matrix block at a neighboring physical pair. \\
$B^s,\widetilde B^s$ & $\chi\times\chi$ & Right-copy and left-copy exact MPS embeddings of $T$. \\
$A^s$ & $D\times D$ & Ordinary MPS block in VUMPS. \\
$l,r$ & $\chi\times\chi$ for LETTA & Dominant left and right fixed points of a normalized structured transfer map. \\
$l_p,r_p$ & $D\times D$ & Conditioned fixed-point blocks used to canonicalize LETTA. \\
$\rho_q$ & $D\times D$ & Right-metric block of the already left-canonical $T_L$; this is used for tangent whitening and is kept distinct from $r_p$. \\
$T_L,T_C,T_R$ & $D\times d\times d\times D$ & Conditional left-canonical, center, and right-canonical LETTA tensors. \\
$C_p$ & $D\times D$ & Physical-sector-conditioned LETTA center matrix. \\
$A_L,A_C,A_R,C$ & VUMPS canonical objects & Ordinary one-site MPS mixed-canonical tensors and its single center matrix. \\
$\langle X,Y\rangle$ & $\Tr(X^\dagger Y)$ & Hilbert--Schmidt inner product; vectorized formulas use the same convention. \\
$\epsilon_{\mathrm{tan}}$ & \code{residual_norm} & Gauge-invariant conditional tangent residual in the default VULETTA path. \\
$\epsilon_{\mathrm{can}}$ &
\code{canonical_}\allowbreak\code{residual_norm} &
Maximum conditional isometry or center-relation error. \\
\bottomrule
\end{longtable}
\endgroup

The rank-four Hamiltonian order in both packages is
\begin{equation}
  h_{s t u v}=\langle s,t|h|u,v\rangle.
  \label{eq:h-index-order}
\end{equation}
A matrix of shape $(d^2,d^2)$ is reshaped into this order without a
permutation.  This convention is important when translating the
\code{einsum} strings into equations.

\section{The uniform nearest-neighbor LETTA state}
\label{sec:vuletta}

\subsection{Finite periodic amplitude}

For a periodic configuration $\bm{s}=(s_1,\ldots,s_N)$ with
$s_{N+1}=s_1$, the implemented amplitude is
\begin{equation}
  \Psi_T(s_1,\ldots,s_N)
  =\Tr\!\left[
    T^{s_1,s_2}T^{s_2,s_3}\cdots T^{s_N,s_1}
  \right].
  \label{eq:letta-amplitude}
\end{equation}
This is exactly the loop in
\code{UniformLETTA.periodic_amplitude}.  Each matrix block uses the physical
state of one site and its right neighbor.  The same $T$ is repeated at every
position, so the ansatz is one-site translation invariant despite its pair
label.

The raw tensor contains
\begin{equation}
  n_T=d^2D^2
  \label{eq:letta-entry-count}
\end{equation}
complex scalar entries, or that many real entries when \code{real=True}.
Frobenius normalization of $T$ changes finite-state amplitudes by a global
length-dependent factor but leaves normalized thermodynamic observables
unchanged.

\subsection{Two exact structured-MPS embeddings}

Define the composite index $(\alpha,p)$ of dimension $\chi=dD$.  The primary
embedding is
\begin{equation}
  B^s_{(\alpha,p),(\beta,q)}
  =\delta_{q,s}\,T^{p,s}_{\alpha\beta}.
  \label{eq:right-copy-embedding}
\end{equation}
Only the slice whose outgoing copied state $q$ equals the physical MPS index
$s$ is nonzero.  Matrix multiplication imposes the next incoming copied state,
and hence
\begin{align}
  \Tr(B^{s_1}\cdots B^{s_N})
  &=\sum_{\substack{\alpha_1,\ldots,\alpha_N\\
                    p_1,\ldots,p_N}}
    \prod_{n=1}^N
    \delta_{p_{n+1},s_n}
    T^{p_n,s_n}_{\alpha_n\alpha_{n+1}}
  \notag\\
  &=\Tr\!\left[
    T^{s_N,s_1}T^{s_1,s_2}\cdots T^{s_{N-1},s_N}
  \right]
  =\Psi_T(\bm{s}),
  \label{eq:right-copy-proof}
\end{align}
where cyclicity of the trace was used in the last equality.

The shifted embedding in \code{shifted_structured_mps_tensor} moves the copy
constraint to the incoming leg,
\begin{equation}
  \widetilde B^s_{(\alpha,p),(\beta,q)}
  =\delta_{p,s}\,T^{p,q}_{\alpha\beta}.
  \label{eq:left-copy-embedding}
\end{equation}
Now multiplication directly enforces $q_n=s_{n+1}$, giving
\begin{equation}
  \Tr(\widetilde B^{s_1}\cdots\widetilde B^{s_N})
  =\Tr(T^{s_1,s_2}\cdots T^{s_N,s_1}).
  \label{eq:left-copy-proof}
\end{equation}
The two embeddings represent the same amplitudes.  They are both required
because one exposes the conditioned left fixed points and the other exposes
the conditioned right fixed points used by the exact gauge construction.

\paragraph{Crucial implementation boundary.}
The arrays $B$ and $\widetilde B$ contain many structural zeros and have bond
$dD$, but the optimizer never treats those entries as independent variables.
Every objective direction is first made in $T$ and then embedded with the same
copy rule.  The variational manifold therefore remains tied at every step.

\section{Thermodynamic transfer calculus for LETTA}
\label{sec:transfer}

\subsection{Transfer maps and normalization}

For either structured tensor $A^s$ (usually $A^s=B^s$ in contractions),
define
\begin{equation}
  \E_A(X)=\sum_s A^sX(A^s)^\dagger,
  \qquad
  \E_A^\dagger(Y)=\sum_s(A^s)^\dagger YA^s.
  \label{eq:transfer-maps}
\end{equation}
The dense matrix built by \code{_dense_transfer_matrix} satisfies
\begin{equation}
  [\bm E_A]_{(a,e),(b,c)}
  =\sum_s A^s_{ab}(A^s_{ec})^*.
  \label{eq:dense-transfer-index}
\end{equation}
It acts on $\vecop(X)$ with matrix dimension $\chi^2\times\chi^2$.

Let $\lambda_0$ be the spectral radius of the unnormalized transfer matrix.
The implementation rejects a numerically zero radius and rejects a
noninjective state when
\begin{equation}
  \Delta_{\E}
  =1-\frac{|\lambda_1|}{|\lambda_0|}
  \leq \epsilon_{\mathrm{inj}},
  \qquad \epsilon_{\mathrm{inj}}=10^{-10}\ \text{by default}.
  \label{eq:injectivity-gap}
\end{equation}
The normalized structured tensor is
\begin{equation}
  A^s=\frac{B^s}{\sqrt{\lambda_0}},
  \label{eq:normalized-structured-tensor}
\end{equation}
so its dominant transfer eigenvalue is one.  Positive Hermitian fixed points
are chosen and normalized as
\begin{equation}
  \E_A(r)=r,
  \qquad
  \E_A^\dagger(l)=l,
  \qquad
  \Tr r=1,
  \qquad
  \Tr(lr)=1.
  \label{eq:fixed-point-normalization}
\end{equation}
The helper \code{_positive_eigenmatrix} phase-aligns, Hermitizes, and checks the
dominant eigenmatrix.  A significantly negative eigenvalue is an error, not
silently clipped into a density matrix.

\subsection{One-site and two-site observables}

For a one-site operator $O_{st}=\langle s|O|t\rangle$, define the insertion
map
\begin{equation}
  \mathcal O_A(X)=\sum_{s,t}O_{st}A^tX(A^s)^\dagger.
  \label{eq:one-site-map}
\end{equation}
The thermodynamic expectation is
\begin{equation}
  \langle O\rangle=(l|\mathcal O_A|r)
  =\sum_{s,t}O_{st}\Tr\!\left[l^\dagger A^t r(A^s)^\dagger\right].
  \label{eq:one-site-expectation}
\end{equation}
This is the contraction in \code{one_site_expectation}.

Introduce the two-site product tensor
\begin{equation}
  P^{st}=A^sA^t,
  \qquad
  P_{a,s,t,c}=\sum_b A^s_{ab}A^t_{bc}.
  \label{eq:pair-tensor}
\end{equation}
For $h_{stuv}=\langle s,t|h|u,v\rangle$, the local Hamiltonian map is
\begin{equation}
  \Htwo_A(X)
  =\sum_{s,t,u,v}h_{stuv}
   A^uA^vX(A^sA^t)^\dagger.
  \label{eq:two-site-map}
\end{equation}
Therefore
\begin{equation}
  e=(l|\Htwo_A|r)
  =\sum_{s,t,u,v}h_{stuv}
   \Tr\!\left[l^\dagger A^uA^v r(A^sA^t)^\dagger\right].
  \label{eq:letta-energy}
\end{equation}
The same formula evaluates any two-site observable.  Hermiticity is checked
before an energy call, and an imaginary energy above $10^{-9}$ is rejected.

\section{Exact physical-sector-dependent LETTA gauge}
\label{sec:conditional-gauge}

\subsection{Gauge group and telescoping invariance}

Choose one invertible virtual matrix $G_p\in\mathrm{GL}(D,\CC)$ for every
physical state $p$.  The LETTA block transformation is
\begin{equation}
  T^{p,q}\longmapsto T'^{p,q}
  =G_p^{-1}T^{p,q}G_q.
  \label{eq:letta-gauge}
\end{equation}
Substitution into Eq.~\eqref{eq:letta-amplitude} gives
\begin{align}
  \Psi_{T'}(\bm{s})
  &=\Tr\!\left[
    G_{s_1}^{-1}T^{s_1,s_2}G_{s_2}
    G_{s_2}^{-1}T^{s_2,s_3}G_{s_3}\cdots
    G_{s_N}^{-1}T^{s_N,s_1}G_{s_1}
  \right]
  \notag\\
  &=\Psi_T(\bm{s}).
  \label{eq:gauge-telescope}
\end{align}
This is implemented by \code{UniformLETTA.gauge_transform}.  It is strictly
larger than the ordinary MPS gauge $A^s\mapsto G^{-1}A^sG$ because $G_p$ may
depend on the shared physical state.

For infinitesimal generators $K_p$, Eq.~\eqref{eq:letta-gauge} gives the null
direction
\begin{equation}
  \delta_K T^{p,q}=-K_pT^{p,q}+T^{p,q}K_q.
  \label{eq:infinitesimal-gauge}
\end{equation}
The analytic gradient tests verify
\begin{equation}
  \RePart\sum_{p,q}\Tr\!\left[(g^{p,q})^\dagger
  \delta_KT^{p,q}\right]=0.
  \label{eq:gradient-gauge-null}
\end{equation}

\subsection{Why the transfer fixed points split into sectors}

Let
\begin{equation}
  T_0^{p,q}=\frac{T^{p,q}}{\sqrt{\lambda_0}}
  \label{eq:T0}
\end{equation}
be normalized by the right-copy transfer radius.  The copy constraints force
the relevant structured fixed points to be block diagonal in the shared
physical index.  Write those $D\times D$ blocks as $l_p$ and $r_p$.  They obey
\begin{equation}
  l_q=\sum_p(T_0^{p,q})^\dagger l_pT_0^{p,q},
  \label{eq:left-sector-fixed}
\end{equation}
from the left fixed point of the right-copy embedding, and
\begin{equation}
  r_p=\sum_qT_0^{p,q}r_q(T_0^{p,q})^\dagger,
  \label{eq:right-sector-fixed}
\end{equation}
from the right fixed point of the shifted embedding.  The code explicitly
checks that off-diagonal physical blocks are negligible before extracting
these sectors.

For each positive semidefinite block $x$, the helper
\code{_conditioned_positive_sqrt} diagonalizes
\begin{equation}
  x=U\diag(\sigma_1,\ldots,\sigma_D)U^\dagger,
  \label{eq:positive-diagonalization}
\end{equation}
and constructs
\begin{align}
  x^{1/2}&=U\diag(\sqrt{\sigma_i})U^\dagger,\\
  x^{-1/2}_{+}&=U\diag\!\left(
    \mathbf 1_{\sigma_i>\tau\sigma_{\max}}\sigma_i^{-1/2}
  \right)U^\dagger,
  \label{eq:pseudoinverse-sqrt}
\end{align}
where $\tau=\code{canonical_rcond}=10^{-12}$ by default.  A negative
eigenvalue below the numerical allowance is rejected.  Active canonicalization
requires every $l_p$ and $r_p$ to have rank $D$ unless
\code{allow_rank_deficient=True} is explicitly requested.

\subsection{Conditional mixed-canonical construction}

The implemented canonical tensors are
\begin{align}
  T_L^{p,q}&=l_p^{1/2}T_0^{p,q}l_q^{-1/2},
  \label{eq:TL}\\
  T_R^{p,q}&=r_p^{-1/2}T_0^{p,q}r_q^{1/2}.
  \label{eq:TR}
\end{align}
Using Eqs.~\eqref{eq:left-sector-fixed} and
\eqref{eq:right-sector-fixed}, one obtains the sectorwise isometries
\begin{align}
  \sum_p(T_L^{p,q})^\dagger T_L^{p,q}&=\Id_D
  &&\text{for every }q,
  \label{eq:left-conditional-isometry}\\
  \sum_qT_R^{p,q}(T_R^{p,q})^\dagger&=\Id_D
  &&\text{for every }p.
  \label{eq:right-conditional-isometry}
\end{align}

Define unnormalized sector centers
\begin{equation}
  \widehat C_p=l_p^{1/2}r_p^{1/2},
  \qquad
  \mathcal N_C=
  \left(\sum_p\norm{\widehat C_p}_F^2\right)^{1/2},
  \qquad
  C_p=\widehat C_p/\mathcal N_C.
  \label{eq:conditional-centers}
\end{equation}
The center tensor is
\begin{equation}
  T_C^{p,q}=T_L^{p,q}C_q.
  \label{eq:TC-definition}
\end{equation}
The two center factorizations agree exactly in full support,
\begin{align}
  T_L^{p,q}C_q
  &=\frac{1}{\mathcal N_C}
    l_p^{1/2}T_0^{p,q}r_q^{1/2}
  \notag\\
  &=C_pT_R^{p,q}.
  \label{eq:center-relation-proof}
\end{align}
Thus
\begin{equation}
  T_C^{p,q}=T_L^{p,q}C_q=C_pT_R^{p,q}.
  \label{eq:center-relation}
\end{equation}
Moreover, Eq.~\eqref{eq:left-conditional-isometry} implies
\begin{equation}
  \norm{T_C}_F^2
  =\sum_q\Tr\!\left[
    C_q^\dagger\left(\sum_p(T_L^{p,q})^\dagger T_L^{p,q}\right)C_q
  \right]
  =\sum_q\norm{C_q}_F^2=1.
  \label{eq:TC-norm}
\end{equation}

The normalized $T_L$ differs from the input by the transfer normalization and
a telescoping gauge.  Consequently a length-$N$ periodic amplitude obeys
\begin{equation}
  \Psi_{T_L}(\bm{s})
  =\lambda_0^{-N/2}\Psi_T(\bm{s}).
  \label{eq:amplitude-scale}
\end{equation}
This factor is returned by \code{ConditionalCanonicalLETTA.amplitude_scale}.
All normalized thermodynamic observables are invariant.

\subsection{Canonical validation and numerical boundary}

The class \code{ConditionalCanonicalLETTA} validates
\begin{equation}
  \epsilon_L=\max_q\norm{
    \sum_p(T_L^{p,q})^\dagger T_L^{p,q}-\Id}_F,
  \label{eq:left-error}
\end{equation}
\begin{equation}
  \epsilon_R=\max_p\norm{
    \sum_qT_R^{p,q}(T_R^{p,q})^\dagger-\Id}_F,
  \label{eq:right-error}
\end{equation}
and
\begin{equation}
  \epsilon_C=\max_{p,q}\left\{
  \norm{T_C^{p,q}-T_L^{p,q}C_q}_F,
  \norm{T_L^{p,q}C_q-C_pT_R^{p,q}}_F
  \right\}.
  \label{eq:center-error}
\end{equation}
The default validation tolerance is $10^{-8}$.  It also checks
$\norm{C}_F=\norm{T_C}_F=1$.  Finally, canonicalization recomputes transfer
data for the returned $T_L$.  This last check matters near an injectivity
boundary: a raw trial and its ill-conditioned gauge transform can straddle the
numerical spectral-gap threshold.  Such a retraction is rejected rather than
passed into the reduced resolvent.

\section{Analytic transfer-response gradient}
\label{sec:gradient}

\subsection{Transfer variations}

Let $A$ denote the normalized structured MPS tensor of
Eq.~\eqref{eq:normalized-structured-tensor}, and let $V$ be a structured
variation induced by a variation of $T$.  The full first-order transfer
variation is
\begin{equation}
  \delta\E_{A,V}(X)
  =\sum_s\left[
    V^sX(A^s)^\dagger+A^sX(V^s)^\dagger
  \right].
  \label{eq:transfer-variation}
\end{equation}
The implementation also separates the ket and bra pieces,
\begin{align}
  \mathcal K_V(X)&=\sum_sV^sX(A^s)^\dagger,
  \label{eq:ket-variation}\\
  \mathcal B_V(X)&=\sum_sA^sX(V^s)^\dagger,
  \label{eq:bra-variation}
\end{align}
so $\delta\E_{A,V}=\mathcal K_V+\mathcal B_V$.  A double insertion used by
the dense tangent metric is
\begin{equation}
  \mathcal D_{U,V}(X)=\sum_sV^sX(U^s)^\dagger.
  \label{eq:double-variation}
\end{equation}
Four private direction helpers implement, respectively, the full transfer
variation, its ket part, its bra part, and the double insertion.

The two-site tensor varies as
\begin{equation}
  \delta P^{st}=V^sA^t+A^sV^t.
  \label{eq:pair-variation}
\end{equation}
Substitution into Eq.~\eqref{eq:two-site-map} produces
\begin{align}
  \delta\Htwo_{A,V}(X)
  =\sum_{s,t,u,v}h_{stuv}\big[&
    \delta P^{uv}X(P^{st})^\dagger
    +P^{uv}X(\delta P^{st})^\dagger
  \big].
  \label{eq:two-site-variation}
\end{align}

\subsection{Removing the dominant transfer mode}

Use vectorized fixed points $|r)$ and $(l|$ with $(l|r)=1$.  Define
\begin{equation}
  \Proj=|r)(l|,
  \qquad
  \Qproj=\Id-\Proj.
  \label{eq:PQ}
\end{equation}
The reduced resolvent is
\begin{equation}
  \Res
  =\Qproj(\Id-\E_A)^{-1}\Qproj
  =(\Id-\bm E_A+\Proj)^{-1}-\Proj.
  \label{eq:reduced-resolvent}
\end{equation}
The second equality is the dense formula implemented in
\code{_reduced_resolvent}.  It removes the singular eigenvalue-one mode and
sums all connected transfer responses.  Injectivity is required so that no
other unit-modulus eigenmode makes this inverse singular.

\subsection{Normalization derivative for a raw LETTA coordinate}

The energy is invariant under an overall nonzero scaling of $T$, but the
structured contraction tensor is explicitly normalized by $\lambda_0(T)$.
For a raw embedded direction $\delta B$, first define
\begin{equation}
  V_{\mathrm{raw}}=\frac{\delta B}{\sqrt{\lambda_0}}.
  \label{eq:raw-normalized-direction}
\end{equation}
The relative spectral-radius change evaluated by the code is
\begin{equation}
  \eta=(l|\delta\E_{A,V_{\mathrm{raw}}}|r)
  =\frac{\delta\lambda_0}{\lambda_0}.
  \label{eq:relative-eigenvalue-change}
\end{equation}
The direction of the normalized structured tensor is therefore
\begin{equation}
  V=V_{\mathrm{raw}}-\frac{1}{2}\RePart(\eta)A.
  \label{eq:spectral-normalization-direction}
\end{equation}
This correction removes the radial transfer-growth direction before any
energy response is evaluated.

\subsection{Directional energy derivative}

Let
\begin{equation}
  U=\Htwo_A(r),
  \qquad
  L_H=\Res U,
  \qquad
  R_V=\Res\,\delta\E_{A,V}(r).
  \label{eq:gradient-environments}
\end{equation}
Differentiating Eq.~\eqref{eq:letta-energy}, including both fixed-point
responses, gives the implemented expression
\begin{align}
  \delta e[V]
  ={}&(l|\delta\Htwo_{A,V}|r)
  +(l|\delta\E_{A,V}\Res\Htwo_A|r)
  +(l|\Htwo_A\Res\delta\E_{A,V}|r)
  \notag\\
  ={}&(l|\delta\Htwo_{A,V}(r))
  +(l|\delta\E_{A,V}(L_H))
  +(l|\Htwo_A(R_V)).
  \label{eq:directional-energy-gradient}
\end{align}
This is \code{_directional_energy_derivative}.  The first term differentiates
the local two-site insertion.  The second and third terms account for the
semi-infinite left and right transfer responses.  The code rejects an
imaginary component above $10^{-8}$ for a Hermitian Hamiltonian.

\subsection{Assembly of the tensor gradient}

For every LETTA scalar coordinate $n$, the code creates a basis direction
$E_n$ in $T$, embeds it using exactly the sparse rule
Eq.~\eqref{eq:right-copy-embedding}, applies
Eq.~\eqref{eq:spectral-normalization-direction}, and evaluates
Eq.~\eqref{eq:directional-energy-gradient}.  For real $T$,
\begin{equation}
  g_n=\delta e[E_n].
  \label{eq:real-gradient-coordinate}
\end{equation}
For complex $T$, it separately evaluates $E_n$ and $iE_n$ and stores
\begin{equation}
  g_n=\delta e[E_n]+i\,\delta e[iE_n].
  \label{eq:complex-gradient-coordinate}
\end{equation}
The convention is
\begin{equation}
  \delta e=\RePart\langle g,\delta T\rangle_F
  =\RePart\sum_{p,q}\Tr\!\left[(g^{p,q})^\dagger\delta T^{p,q}\right].
  \label{eq:complex-gradient-convention}
\end{equation}
Central-difference tests cover real coordinates and both real and imaginary
directions of complex coordinates.

\paragraph{Meaning of an analytic gradient in this code.}
No energy finite difference is used in \code{energy_and_gradient}.  The code
does loop over the $d^2D^2$ independent tied tensor entries, but every
coordinate derivative is computed from the same transfer eigensystem and the
same reduced resolvent.  The implementation is analytic but not yet organized
as a single reverse-mode tensor contraction.

\section{Gauge-fixed conditional tangent space}
\label{sec:tangent}

\subsection{Left-horizontal condition}

For each current physical state $q$, stack the left-canonical blocks into
\begin{equation}
  M_q=
  \begin{bmatrix}
    T_L^{0,q}\\T_L^{1,q}\\\vdots\\T_L^{d-1,q}
  \end{bmatrix}
  \in\CC^{dD\times D}.
  \label{eq:Mq}
\end{equation}
Equation~\eqref{eq:left-conditional-isometry} is simply
\begin{equation}
  M_q^\dagger M_q=\Id_D.
  \label{eq:Mq-isometry}
\end{equation}
An active variation $X^{p,q}$ must satisfy
\begin{equation}
  \sum_p(T_L^{p,q})^\dagger X^{p,q}=0
  \quad\Longleftrightarrow\quad
  M_q^\dagger X_q=0
  \qquad\text{for each }q,
  \label{eq:horizontal-condition}
\end{equation}
where $X_q$ stacks $X^{p,q}$ in the same way as $M_q$.  This is the conditional
analogue of the usual left-gauge tangent condition for an MPS.

A complete QR factorization produces
\begin{equation}
  [M_q\;N_q]\ \text{unitary},
  \qquad
  N_q^\dagger M_q=0,
  \qquad
  N_q\in\CC^{dD\times(d-1)D}.
  \label{eq:Nq}
\end{equation}
All horizontal variations can be written as $X_q=N_qY_q$.

\subsection{Right metric and whitened coordinates}

Recompute the right transfer fixed point of the structured MPS represented by
$T_L$, reshape it as
\begin{equation}
  r_{(\alpha,q),(\beta,q')},
  \label{eq:TL-right-fixed}
\end{equation}
and extract its diagonal blocks
\begin{equation}
  \rho_q=r_{(\alpha,q),(\beta,q)}.
  \label{eq:rhoq}
\end{equation}
These $\rho_q$ are tangent-metric blocks.  They are not merely renamed copies
of the shifted-embedding $r_p$ in the canonicalization derivation.

The implementation whitens the horizontal coordinate on the right,
\begin{equation}
  X_q=N_qZ_q\rho_q^{-1/2,+},
  \qquad
  Z_q\in\CC^{(d-1)D\times D}.
  \label{eq:whitened-tangent}
\end{equation}
The superscript $+$ denotes the numerical pseudoinverse square root with the
same \code{canonical_}\allowbreak\code{rcond} threshold.  If $\rho_q$ loses
numerical rank, only its supported eigenspace remains active.

\subsection{Gradient pullback, descent, and residual}

Stack the tensor gradient at fixed $q$ into $g_q\in\CC^{dD\times D}$.  Pulling
Eq.~\eqref{eq:complex-gradient-convention} back through
Eq.~\eqref{eq:whitened-tangent} gives
\begin{equation}
  g_{Z_q}=N_q^\dagger g_q\rho_q^{-1/2,+}.
  \label{eq:reduced-gradient}
\end{equation}
Steepest descent in the whitened coordinate chooses
\begin{equation}
  \delta Z_q=-g_{Z_q},
  \qquad
  X_q=-N_qg_{Z_q}\rho_q^{-1/2,+}.
  \label{eq:conditional-descent}
\end{equation}
The energy slope is exactly
\begin{align}
  \RePart\langle g,X\rangle_F
  &= -\sum_q\Tr(g_{Z_q}^\dagger g_{Z_q})
  \notag\\
  &=-\epsilon_{\mathrm{tan}}^2,
  \label{eq:slope-residual}
\end{align}
where the reported residual is
\begin{equation}
  \epsilon_{\mathrm{tan}}
  =\left(\sum_q\norm{g_{Z_q}}_F^2\right)^{1/2}.
  \label{eq:tangent-residual}
\end{equation}
This identity is directly tested for real and complex states.

If every $\rho_q$ has full rank, the complex tangent dimension is
\begin{equation}
  n_{\mathrm{tan}}^{\CC}=d(d-1)D^2.
  \label{eq:full-complex-tangent-dim}
\end{equation}
For a real calculation this is also the real coordinate count.  For a complex
calculation represented by real and imaginary parts, the reported real count is
\begin{equation}
  n_{\mathrm{tan}}^{\R}=2d(d-1)D^2.
  \label{eq:full-real-tangent-dim}
\end{equation}
With numerical sector ranks $r_q=\rank(\rho_q)$, the implementation reports
\begin{equation}
  n_{\mathrm{supported}}^{\CC}
  =\sum_q(d-1)D\,r_q,
  \label{eq:supported-tangent-dim}
\end{equation}
and doubles this value for complex real-packed coordinates.

\subsection{Why this removes scale, phase, and gauge}

The raw complex parameter count is $d^2D^2$.  A generic projective LETTA state
has the sector-dependent gauge of Eq.~\eqref{eq:letta-gauge}, plus overall
complex scaling.  The gauge generators contain $dD^2$ complex entries, but
the common scalar generator is trivial; adding projective scale restores a
total redundancy of $dD^2$.  Thus the expected generic quotient dimension is
\begin{equation}
  d^2D^2-dD^2=d(d-1)D^2,
  \label{eq:dimension-count-quotient}
\end{equation}
matching Eq.~\eqref{eq:full-complex-tangent-dim}.  The horizontal complement
implements this quotient without explicitly building or inverting all gauge
generators.

\section{Dense connected tangent metric: legacy path and oracle}
\label{sec:dense-gram}

The implementation retains a dense connected Gram matrix for validation and
the legacy natural-gradient update.  It is not called by the default
conditional-canonical solver.

Let $V_\mu$ be normalized structured parameter directions, including
Frobenius normalization of raw $T$ and the spectral correction in
Eq.~\eqref{eq:spectral-normalization-direction}.  Define
\begin{align}
  K_\nu&=\mathcal K_{V_\nu}(r),
  &B_\mu&=\mathcal B_{V_\mu}(r),
  \label{eq:KB-sources}\\
  k_\nu&=(l|K_\nu),
  &b_\mu&=(l|B_\mu).
  \label{eq:KB-parallel}
\end{align}
The unsymmetrized connected overlap used by the code is
\begin{align}
  \widetilde G_{\mu\nu}
  ={}&(l|\mathcal D_{V_\mu,V_\nu}(r))
  -b_\mu k_\nu
  \notag\\
  &+(l|\mathcal B_{V_\mu}(\Res K_\nu))
  +(l|\mathcal K_{V_\nu}(\Res B_\mu)).
  \label{eq:connected-gram}
\end{align}
The returned real metric is
\begin{equation}
  G=\frac{1}{2}\left[
    \RePart\widetilde G+(\RePart\widetilde G)^T
  \right].
  \label{eq:symmetrized-gram}
\end{equation}
Radial, phase, and virtual-gauge variations appear as numerical null modes.

For a packed real coordinate gradient $g$, \code{natural_gradient} diagonalizes
\begin{equation}
  G=U\diag(\gamma_i)U^T
  \label{eq:metric-eigendecomposition}
\end{equation}
and retains $\gamma_i>\max(\tau\gamma_{\max},\epsilon_{\mathrm{mach}})$.  It
returns
\begin{equation}
  \eta=G^+g,
  \qquad
  \epsilon_G=\sqrt{g^TG^+g},
  \qquad
  r_G=\rank_{\tau}(G).
  \label{eq:dense-natural-gradient}
\end{equation}
The legacy solver uses $-\eta$ as its descent.  For small systems the tests
verify that
\begin{equation}
  \epsilon_G=\epsilon_{\mathrm{tan}}
  \label{eq:oracle-equality}
\end{equation}
and that the rank agrees with the conditional coordinate count.  This is the
main numerical evidence that the direct conditional construction fully uses
the LETTA gauge rather than merely imposing parameter normalization.

\section{The default VULETTA optimization algorithm}
\label{sec:vuletta-solver}

\subsection{Inputs, normalization, and the three update paths}

The public entry point is \code{solver.vuletta}.  It accepts a Hermitian
nearest-neighbor operator either as a matrix of shape $d^2\times d^2$ or as
$h_{stuv}$ of shape $d\times d\times d\times d$.  The convention is
\begin{equation}
  h_{stuv}=\langle s,t|\widehat h|u,v\rangle.
  \label{eq:h-index-convention}
\end{equation}
The current optimizer has three update methods:
\begin{enumerate}
  \item \code{conditional_canonical}, the default and the method whose gauge
        geometry is derived in Sections~\ref{sec:conditional-gauge} and
        \ref{sec:tangent};
  \item \code{natural_gradient}, the dense connected-metric implementation of
        Section~\ref{sec:dense-gram};
  \item \code{lbfgs}, an ordinary Euclidean L-BFGS-B optimization of normalized
        tensor coordinates, kept as a baseline.
\end{enumerate}
The first two paths require the analytic gradient.  The L-BFGS path may use
either that gradient or a two- or three-point finite difference.

The helper \code{_unpack_tensor} normalizes every raw L-BFGS parameter vector:
\begin{equation}
  T(x)=\frac{\operatorname{unpack}(x)}
             {\norm{\operatorname{unpack}(x)}_F}.
  \label{eq:unpack-normalization}
\end{equation}
For complex optimization, real and imaginary parts are packed consecutively.
For \code{real=True}, only real coordinates are retained and a complex initial
tensor is rejected.

\subsection{One conditional-canonical iteration}

Let $T_k$ already be represented by a valid
\code{ConditionalCanonicalLETTA}.  One default iteration performs the
following operations.

\paragraph{1. Evaluate the physical objective and its tensor gradient.}
\begin{equation}
  (e_k,g_k)=\operatorname{energy\_and\_gradient}(T_k,h),
  \qquad
  \delta e_k=\RePart\langle g_k,\delta T\rangle_F .
  \label{eq:solver-energy-gradient}
\end{equation}

\paragraph{2. Pull the gradient to horizontal conditional coordinates.}
For each output sector $q$, the construction of
Eq.~\eqref{eq:reduced-gradient} gives $g_{Z_q}$.  The reported tangent
residual is
\begin{equation}
  \epsilon_k
  =
  \left(\sum_q\norm{g_{Z_q}}_F^2\right)^{1/2}.
  \label{eq:solver-tangent-residual}
\end{equation}
The unscaled descent $X_k$ is Eq.~\eqref{eq:conditional-descent}, and hence
\begin{equation}
  m_k
  \equiv
  \RePart\langle g_k,X_k\rangle_F
  =-\epsilon_k^2<0
  \label{eq:exact-descent-slope}
\end{equation}
unless the state is stationary on the supported tangent space.

\paragraph{3. Check both kinds of stationarity.}
The canonical residual stored by the solver is
\begin{equation}
  c_k
  =
  \max\left\{
    \epsilon_L,\epsilon_R,\epsilon_C
  \right\},
  \label{eq:combined-canonical-residual}
\end{equation}
where the three terms are the left-isometry, right-isometry, and center
compatibility errors of
Eqs.~\eqref{eq:left-conditional-isometry},
\eqref{eq:right-conditional-isometry}, and
\eqref{eq:center-relation}.  Convergence is declared only when
\begin{equation}
  \epsilon_k\leq\tau_{\mathrm{stat}}
  \quad\text{and}\quad
  c_k\leq\tau_{\mathrm{stat}}.
  \label{eq:two-part-convergence}
\end{equation}
This is important: a small physical gradient does not excuse a badly
canonicalized representative, and a perfect gauge does not imply an energy
stationary point.

\paragraph{4. Cap the whitened-coordinate step.}
The code uses $\epsilon_k$ as the coordinate norm of the natural descent.  If
$\epsilon_k>\Delta_{\max}$ it rescales
\begin{equation}
  X_k\leftarrow
  \frac{\Delta_{\max}}{\epsilon_k}X_k.
  \label{eq:step-cap}
\end{equation}
The directional slope is then recomputed from the actual scaled $X_k$.  A
nonfinite or nonnegative slope terminates the iteration rather than sending
an invalid direction to line search.

\paragraph{5. Armijo line search with canonical retraction.}
Starting at $a=1$, the trial tensor is
\begin{equation}
  T_{\mathrm{raw}}(a)=T_k+aX_k.
  \label{eq:raw-trial-tensor}
\end{equation}
It first has to satisfy
\begin{equation}
  e\!\left(T_{\mathrm{raw}}(a)\right)
  \leq e_k+c_{\mathrm A}a\,m_k,
  \qquad 0<c_{\mathrm A}<1.
  \label{eq:armijo}
\end{equation}
The trial is then passed through \code{conditional_canonicalize}.  Thus
canonicalization is the retraction back to the chosen gauge section:
\begin{equation}
  T_{k+1}
  =
  \operatorname{Canon}\!\left[T_k+aX_k\right].
  \label{eq:conditional-retraction}
\end{equation}
If the energy evaluation, injectivity test, matrix square root, rank test, or
canonical validation fails, the trial is rejected.  The step is halved,
$a\leftarrow a/2$, until accepted or until the line-search/evaluation budget
is exhausted.  Retraction is therefore part of the acceptance rule rather
than a cosmetic postprocessing operation.

\subsection{Compact pseudocode}

\begin{enumerate}
  \item Build or validate $T_0$, normalize its Frobenius norm, and compute
        $\mathcal C_0=\operatorname{Canon}(T_0)$.
  \item For $k=0,\ldots,K_{\max}-1$:
    \begin{enumerate}
      \item evaluate $(e_k,g_k)$;
      \item construct the conditional horizontal tangent, residual
            $\epsilon_k$, and dimension;
      \item compute $c_k$ and stop if Eq.~\eqref{eq:two-part-convergence}
            holds;
      \item cap $X_k$ by Eq.~\eqref{eq:step-cap};
      \item backtrack until Eqs.~\eqref{eq:armijo} and
            \eqref{eq:conditional-retraction} both succeed;
      \item append energy, energy change, residuals, and step size to the
            immutable iteration record.
    \end{enumerate}
  \item Re-evaluate the final analytic gradient and conditional tangent.
        The final convergence flag is based on these fresh residuals, not
        merely on the reason the loop exited.
\end{enumerate}

\subsection{Default numerical parameters}

\begin{table}[htbp]
  \centering
  \caption{Defaults in \code{VULETTAOptions}.}
  \label{tab:vuletta-options}
  \begin{tabular}{@{}lll@{}}
    \toprule
    Option & Default & Meaning \\
    \midrule
    \code{max_iterations} & $200$ & accepted outer iterations \\
    \code{max_function_evaluations} & automatic & total energy/gradient budget \\
    \code{tolerance} & $10^{-8}$ & L-BFGS gradient tolerance \\
    \code{energy_tolerance} & $10^{-13}$ & L-BFGS relative objective tolerance \\
    \code{update_method} & \code{conditional_canonical} & update geometry \\
    \code{gradient_method} & \code{analytic} & derivative implementation \\
    \code{finite_difference_scheme} & \code{3-point} & optional L-BFGS stencil \\
    \code{max_line_search_steps} & $30$ & backtracking budget \\
    \code{max_parameter_step} & $2$ & whitened-coordinate trust cap \\
    \code{metric_rcond} & $10^{-10}$ & dense metric cutoff \\
    \code{canonical_rcond} & $10^{-12}$ & sector inverse/root cutoff \\
    \code{armijo_coefficient} & $10^{-4}$ & $c_{\mathrm A}$ in Eq.~\eqref{eq:armijo} \\
    \code{stationarity_tolerance} & $10^{-6}$ & physical and gauge residual target \\
    \code{verbosity} & $0$ & progress printing \\
    \bottomrule
  \end{tabular}
\end{table}

\section{Source-code map: where each equation lives}
\label{sec:source-map}

\begingroup
\small
\begin{longtable}{@{}p{0.24\linewidth}p{0.30\linewidth}p{0.38\linewidth}@{}}
  \caption{Implementation map for the formulas in this note.}
  \label{tab:source-map}\\
  \toprule
  File & Main object or function & Mathematical responsibility \\
  \midrule
  \endfirsthead
  \toprule
  File & Main object or function & Mathematical responsibility \\
  \midrule
  \endhead
  \path{_vuletta/state.py}
    & \code{UniformLETTA}
    & validates $T_{\alpha p q\beta}$, constructs the two exact structured
      MPS embeddings, and evaluates finite periodic amplitudes \\
  \path{_vuletta/state.py}
    & \code{ConditionalCanonical}\allowbreak\code{LETTA}
    & stores $T_L,T_R,T_C,\{C_p\}$ and evaluates isometry/center residuals \\
  \path{_vuletta/state.py}
    & \code{conditional_canonicalize}
    & transfer normalization, injectivity check, sector fixed points, positive
      roots/inverses, Eqs.~\eqref{eq:TL}--%
      \eqref{eq:conditional-centers}, and final validation \\
  \path{_vuletta/state.py}
    & \code{expand_uniform_letta}
    & bond continuation in Eq.~\eqref{eq:bond-continuation} \\
  \path{_vuletta/operators.py}
    & \code{transfer_data}
    & dense transfer eigensystem, spectral gap, positive fixed points, and
      normalization in Eqs.~\eqref{eq:dense-transfer-index}--%
      \eqref{eq:fixed-point-normalization} \\
  \path{_vuletta/operators.py}
    & observable functions
    & one- and two-site contractions in
      Eqs.~\eqref{eq:one-site-expectation} and
      \eqref{eq:letta-energy} \\
  \path{_vuletta/gradients.py}
    & transfer variation helpers
    & $\mathcal K_V,\mathcal B_V,\mathcal D_{V,W}$ and the normalized
      direction $V$ \\
  \path{_vuletta/gradients.py}
    & \code{energy_and_gradient}
    & reduced resolvent and
      Eq.~\eqref{eq:directional-energy-gradient}, assembled by real/imaginary
      basis directions \\
  \path{_vuletta/gradients.py}
    & \code{conditional_tangent_}\allowbreak\code{direction}
    & QR complements, supported right metric, whitened $Z_q$, horizontal
      descent, residual, and dimension \\
  \path{_vuletta/gradients.py}
    & \code{tangent_gram_matrix}, \code{natural_gradient}
    & dense connected metric and pseudoinverse oracle in
      Eqs.~\eqref{eq:connected-gram}--\eqref{eq:dense-natural-gradient} \\
  \path{_vuletta/solver.py}
    & \code{vuletta}
    & the three update paths, convergence contracts, Armijo backtracking,
      canonical retraction, histories, and result diagnostics \\
  \path{_vuletta/benchmarks/one_dimensional_models.py}
    & benchmark driver
    & common models, references, continuation, repeated timing, observables,
      Markdown, CSV, and JSON serialization \\
  \path{_vumps/state.py}
    & \code{CanonicalMPS}, \code{canonicalize}
    & ordinary mixed-canonical equations and fixed-point roots \\
  \path{_vumps/operators.py}
    & environment/effective operators
    & Eqs.~\eqref{eq:vumps-environment-equation},
      \eqref{eq:vumps-effective-center}, and
      \eqref{eq:vumps-effective-bond} \\
  \path{_vumps/solver.py}
    & \code{vumps}
    & lowest eigenpairs, phase alignment, polar gauge matching, and
      fixed-point residuals \\
  \path{tests/test_vuletta.py}
    & focused VULETTA tests
    & embeddings, gauge, canonicalization, gradients, tangent oracle,
      retraction safety, and solver behavior \\
  \path{tests/test_vumps.py}
    & focused VUMPS tests
    & tensor order, environments, covariance, effective operators, canonical
      validation, and matrix-free paths \\
  \bottomrule
\end{longtable}
\endgroup

\subsection{Array axes and mathematical indices}

The most common source of implementation mistakes is silently permuting
physical and virtual axes.  The exact dictionary is
\begin{align}
  \code{T[a, p, q, b]}
  &\longleftrightarrow T^{p,q}_{\alpha\beta},
  &&(a,p,q,b)=(\alpha,p,q,\beta),
  \label{eq:axis-map-T}\\
  \code{A[a, s, b]}
  &\longleftrightarrow A^s_{\alpha\beta},
  &&(a,s,b)=(\alpha,s,\beta).
  \label{eq:axis-map-A}
\end{align}
For the right-copy embedding, the flattened virtual indices are
\begin{equation}
  i=\alpha d+p,
  \qquad
  j=\beta d+q,
  \label{eq:flattened-copy-index}
\end{equation}
and the assignment is exactly
\begin{equation}
  B[i,s,j]
  =
  \delta_{q,s}T[a,p,s,b].
  \label{eq:array-right-copy}
\end{equation}
The dense transfer builder's contraction can be written
\begin{equation}
  (\mathbf E_A)_{(a,d),(b,c)}
  =
  \sum_s A_{a s b}\overline{A_{d s c}},
  \label{eq:dense-transfer-indices}
\end{equation}
so that vectorizing $X_{bc}$ produces
\begin{equation}
  (\mathbf E_A\vecop X)_{(a,d)}
  =
  \left(\sum_sA^sX A^{s\dagger}\right)_{ad}.
  \label{eq:dense-transfer-vectorization}
\end{equation}
This fixes both the conjugation and row-major reshape conventions.

\section{Numerical assumptions, limitations, and failure meanings}
\label{sec:limitations}

\begin{enumerate}
  \item \textbf{Injectivity is required.}
        The dominant transfer eigenvalue must be unique, and the required
        fixed-point blocks must have supported positive rank.  Cat states,
        exactly symmetry-broken noninjective tensors, and strictly padded
        bond enlargements are rejected.
  \item \textbf{The unit cell is one site.}
        Both solvers documented here assume one-site translation invariance.
        A physical state requiring a two-site pattern may appear as a
        fixed-point cycle rather than a solution.
  \item \textbf{The Hamiltonian is nearest neighbor.}
        All energy and gradient formulas use a two-site $h$.  Longer-range
        terms need either an MPO formulation or additional insertion maps.
  \item \textbf{The canonical representative is not globally unique.}
        Residual sector-unitary rotations can change $T_L,T_R,C_p$ while
        leaving the state invariant.  Energies, observables, amplitudes, and
        tangent norms should be compared, not raw canonical entries.
  \item \textbf{Rank cutoffs define the supported geometry.}
        Near rank loss, changing
        \code{canonical_}\allowbreak\code{rcond} can change the active
        coordinate count in Eq.~\eqref{eq:supported-tangent-dim}.
        A small reported residual then means stationarity on the retained
        support, not along discarded singular directions.
  \item \textbf{VULETTA remains a local optimizer.}
        Armijo descent makes accepted steps locally safe but does not remove
        local minima or guarantee the global ground state.  Multiple seeds
        and bond continuation remain useful.
  \item \textbf{Dense transfer scaling is currently severe.}
        Equations~\eqref{eq:dense-transfer-memory} and
        \eqref{eq:dense-transfer-time} limit practical $D$ well before the
        structured parameter count itself becomes large.
  \item \textbf{Real optimization is a restricted manifold.}
        With \code{real=True}, imaginary tangent directions are omitted.
        This is appropriate only when a real representative of the desired
        solution is expected.
  \item \textbf{A nonconverged result is still labeled nonconverged.}
        Small energy error is not substituted for the tangent/canonical
        contract.  Conversely, a stationary point can be physically poor, as
        the $D=1$ Heisenberg LETTA row demonstrates.
  \item \textbf{The present VUMPS comparison has a known Heisenberg failure.}
        The nonconverged rows in Table~\ref{tab:heisenberg-benchmark} must not
        be used to rank ansatz quality.
\end{enumerate}

\section{Practical reading and debugging checklist}
\label{sec:checklist}

When a VULETTA calculation behaves unexpectedly, inspect quantities in this
order:
\begin{enumerate}
  \item Confirm $h_{stuv}$ follows Eq.~\eqref{eq:h-index-convention} and is
        Hermitian after flattening.
  \item Check the transfer spectral gap and fixed-point positivity before
        interpreting any observable.
  \item Verify $\epsilon_L,\epsilon_R,\epsilon_C$ separately; their maximum
        alone does not identify which canonical identity failed.
  \item Compare the analytic gradient with a few random directional finite
        differences using Eq.~\eqref{eq:complex-gradient-convention}.
  \item Confirm $\RePart\langle g,X\rangle_F<0$ and compare it with the
        predicted $-\epsilon_{\mathrm{tan}}^2$ before step capping.
  \item Inspect the accepted Armijo step and whether rejected trials failed
        energy decrease or canonical/injective retraction.
  \item At small $d,D$, compare the conditional residual and dimension with
        the dense Gram oracle.
  \item Only after these algebraic checks compare energies and observables
        against model references.
\end{enumerate}

\section{Derivation of gauge annihilation by the exact gradient}
\label{app:gauge-annihilation}

For infinitesimal generators $K_p$, the vertical direction is
\begin{equation}
  (\delta_KT)^{p,q}
  =
  -K_pT^{p,q}+T^{p,q}K_q.
  \label{eq:vertical-direction-appendix}
\end{equation}
For a finite periodic chain, differentiating the amplitude gives
\begin{align}
  \delta_K\Psi_T
  &=
  \sum_{n=1}^N
  \Tr\!\left[
    T^{s_1,s_2}\cdots
    \left(
      -K_{s_n}T^{s_n,s_{n+1}}
      +T^{s_n,s_{n+1}}K_{s_{n+1}}
    \right)
    \cdots T^{s_N,s_1}
  \right].
  \label{eq:gauge-amplitude-derivative}
\end{align}
The positive $K_{s_{n+1}}$ term at site $n$ cancels the negative
$K_{s_{n+1}}$ term at site $n+1$ after cyclic rearrangement of the trace.
Therefore
\begin{equation}
  \delta_K\Psi_T=0,
  \qquad
  \delta_Ke=0,
  \qquad
  \RePart\langle g,\delta_KT\rangle_F=0.
  \label{eq:gauge-gradient-annihilation}
\end{equation}
The finite-difference and analytic-gradient tests check the last identity
numerically.  The horizontal construction enforces it geometrically by never
including $\delta_KT$ among the $Z_q$ coordinates.

\section{Why the reduced resolvent has the regularized-inverse form}
\label{app:reduced-resolvent}

With the normalization $(l|r)=1$, define
\begin{equation}
  \Proj=|r)(l|,
  \qquad
  \Qproj=I-\Proj.
  \label{eq:projectors-appendix}
\end{equation}
Because $\E\Proj=\Proj\E=\Proj$, the operator
\begin{equation}
  I-\E+\Proj
  \label{eq:regularized-transfer-operator}
\end{equation}
has eigenvalue $1$ on the dominant subspace and agrees with $I-\E$ on the
connected subspace.  Hence
\begin{align}
  \left(I-\E+\Proj\right)^{-1}
  &=\Proj+\Qproj(I-\E)^{-1}\Qproj,
  \label{eq:regularized-inverse-split}\\
  \Res
  &=
  \left(I-\E+\Proj\right)^{-1}-\Proj
  =
  \Qproj(I-\E)^{-1}\Qproj.
  \label{eq:resolvent-equivalence-appendix}
\end{align}
It follows immediately that
\begin{equation}
  \Res|r)=0,
  \qquad
  (l|\Res=0,
  \qquad
  (I-\E)\Res=\Res(I-\E)=\Qproj.
  \label{eq:resolvent-identities}
\end{equation}
These identities explain both the connected environment terms in
Eq.~\eqref{eq:directional-energy-gradient} and the subtraction terms in the
dense Gram matrix.

\section{Commands to reproduce the focused checks and benchmark}
\label{app:commands}

From the repository root:
\begin{lstlisting}[style=pyqedpython,language=bash]
OPENBLAS_NUM_THREADS=1 OMP_NUM_THREADS=1 \
VECLIB_MAXIMUM_THREADS=1 NUMEXPR_NUM_THREADS=1 \
PYTHONPATH=. python -m pytest -q \
  tests/test_vuletta.py \
  tests/test_vuletta_example.py \
  tests/test_vuletta_benchmark.py \
  tests/test_vumps.py \
  tests/test_vumps_example.py
\end{lstlisting}

The full deterministic benchmark command used for the checked-in snapshot is
\begin{lstlisting}[style=pyqedpython,language=bash]
OPENBLAS_NUM_THREADS=1 OMP_NUM_THREADS=1 \
VECLIB_MAXIMUM_THREADS=1 NUMEXPR_NUM_THREADS=1 \
PYTHONPATH=. python \
  -m pyqed._vuletta.benchmarks.one_dimensional_models \
  --letta-bond-dimensions 1 2 3 \
  --mps-bond-dimensions 1 2 3 4 6 \
  --tolerance 1e-7 --max-iterations 300 --repeats 3 \
  --csv pyqed/_vuletta/benchmarks/results/one_dimensional_models.csv \
  --markdown pyqed/_vuletta/benchmarks/results/one_dimensional_models.md
\end{lstlisting}

\section{Compact VULETTA implementation summary}

The current nearest-neighbor VULETTA implementation contracts an exact sparse
MPS embedding of a leg-tied tensor, fixes the full block-diagonal
physical-sector gauge inherited from that embedding, differentiates the
thermodynamic energy through a connected transfer resolvent, removes all
vertical gauge directions with sector QR complements, whitens the remaining
horizontal directions with the supported right fixed-point metric, and
retracts Armijo steps through injective conditional canonicalization; VUMPS
uses the same canonical-MPS principles without the tying constraint, but
updates its unrestricted center tensors by effective eigenproblems.

\section{Remaining solver diagnostics and continuation}

If no explicit evaluation budget is supplied, a conditional or dense
natural-gradient run receives
\begin{equation}
  N_{\mathrm{eval,max}}
  =
  1+K_{\max}(1+N_{\mathrm{ls}}).
  \label{eq:default-evaluation-budget}
\end{equation}

\subsection{How to interpret the result object}

The most invariant stopping diagnostic is \code{residual_norm}, equal to
$\epsilon_{\mathrm{tan}}$ for the default method.  The fields
\code{coordinate_gradient_norm} and \code{gradient_norm} are Euclidean
coordinate diagnostics after radial projection; they are useful for
debugging but are not invariant under the full sector gauge.
\code{reduced_dimension} and \code{metric_rank} report the active conditional
coordinate count on the default path.  The returned
\code{canonical_state} contains $T_L,T_R,T_C,\{C_p\}$, while
\code{state} is the optimized raw-$T$ representation reconstructed from
that canonical object.

\subsection{Bond-dimension continuation}

The convenience routine \code{continue_bond_dimension} embeds an optimized
$D$ tensor into a larger $D'$ tensor by placing it in the leading virtual
block and adding small random entries elsewhere.  With the present default
noise fraction $\zeta=0.03$, the construction has the schematic form
\begin{equation}
  T'=
  \operatorname{normalize}
  \left[
    \begin{pmatrix}
      T&0\\0&0
    \end{pmatrix}
    +\zeta\,\Xi
  \right].
  \label{eq:bond-continuation}
\end{equation}
The noise is essential because a strictly zero new virtual block would be
rank deficient and would fail the injective conditional canonicalization.

\section{Models, reference formulas, and benchmark evidence}
\label{sec:benchmarks}

\subsection{Transverse-field Ising model}

The benchmark convention is
\begin{equation}
  H_{\mathrm{TFIM}}
  =
  -J\sum_n Z_nZ_{n+1}
  -g\sum_n X_n,
  \label{eq:tfim-hamiltonian}
\end{equation}
implemented by the two-site density
\begin{equation}
  h_{\mathrm{TFIM}}
  =
  -J\,Z\otimes Z
  -\frac g2\left(X\otimes\Id+\Id\otimes X\right).
  \label{eq:tfim-bond-hamiltonian}
\end{equation}
With
\begin{equation}
  \varepsilon(k)
  =
  \sqrt{J^2+g^2-2Jg\cos k},
  \label{eq:tfim-dispersion-root}
\end{equation}
the exact thermodynamic references used by the example are
\begin{align}
  e_0
  &=
  -\frac{1}{2\pi}
  \int_0^{2\pi}\varepsilon(k)\,dk,
  \label{eq:tfim-exact-energy}\\
  \langle X\rangle
  &=
  \frac{1}{2\pi}
  \int_0^{2\pi}
  \frac{g-J\cos k}{\varepsilon(k)}\,dk,
  \label{eq:tfim-exact-x}\\
  \langle Z_nZ_{n+1}\rangle
  &=
  \frac{1}{2\pi}
  \int_0^{2\pi}
  \frac{J-g\cos k}{\varepsilon(k)}\,dk.
  \label{eq:tfim-exact-zz}
\end{align}
The implementation evaluates the last integral in a paired
$k,\pi-k$ form to avoid subtractive cancellation when $J\ll g$.

\subsection{Antiferromagnetic spin-1/2 Heisenberg chain}

The second model is
\begin{equation}
  H_{\mathrm H}
  =
  \sum_n\left(
    S_n^xS_{n+1}^x+
    S_n^yS_{n+1}^y+
    S_n^zS_{n+1}^z
  \right),
  \qquad S^\alpha=\frac12\sigma^\alpha.
  \label{eq:heisenberg-hamiltonian}
\end{equation}
The exact energy reference and symmetry-derived observables used by the
benchmark are
\begin{equation}
  e_0=\frac14-\log 2,
  \qquad
  \langle S^z\rangle=0,
  \qquad
  \langle S_n^zS_{n+1}^z\rangle=\frac{e_0}{3}.
  \label{eq:heisenberg-references}
\end{equation}

\subsection{Recorded deterministic benchmark snapshot}

The checked-in snapshot uses $J=1$, $g=1.5$, seed $3$, tolerance $10^{-7}$,
at most $300$ iterations, and the median of three single-threaded timed
solves after one warm-up.  LETTA dimensions are continued from the preceding
$D$ using $3\%$ noise.  The following tables reproduce the essential row
data.  Conv is the solver's own residual-based flag; a good energy marked no
is deliberately not relabeled as converged.

\begin{table}[H]
  \centering
  \caption{TFIM snapshot at $J=1,g=1.5$.  Runtime is seconds.}
  \label{tab:tfim-benchmark}
  \scriptsize
  \begin{tabular}{@{}llrcrrrrr@{}}
    \toprule
    Method & $D$ & $\chi$ & Conv & $e$ & $|e-e_0|$ &
      $\langle X\rangle$ & $\langle ZZ\rangle$ & $(\epsilon,\ t)$ \\
    \midrule
    VULETTA & 1 & 2 & yes & $-1.666666667$ & $5.26\,10^{-3}$ &
      $0.8888889$ & $0.3333333$ & $(5.82\,10^{-8},0.147)$ \\
    VULETTA & 2 & 4 & yes & $-1.671919277$ & $6.94\,10^{-6}$ &
      $0.8773681$ & $0.3558671$ & $(9.48\,10^{-8},0.549)$ \\
    VULETTA & 3 & 6 & no & $-1.671926188$ & $3.34\,10^{-8}$ &
      $0.8773285$ & $0.3559334$ & $(5.57\,10^{-7},1.365)$ \\
    \midrule
    VUMPS & 1 & 1 & yes & $-1.562500000$ & $1.09\,10^{-1}$ &
      $0.7500001$ & $0.4374998$ & $(7.72\,10^{-8},0.022)$ \\
    VUMPS & 2 & 2 & yes & $-1.671736624$ & $1.90\,10^{-4}$ &
      $0.8780762$ & $0.3546223$ & $(4.40\,10^{-8},0.028)$ \\
    VUMPS & 3 & 3 & yes & $-1.671909431$ & $1.68\,10^{-5}$ &
      $0.8774289$ & $0.3557661$ & $(7.52\,10^{-8},0.051)$ \\
    VUMPS & 4 & 4 & yes & $-1.671925967$ & $2.55\,10^{-7}$ &
      $0.8773301$ & $0.3559308$ & $(4.16\,10^{-8},0.057)$ \\
    VUMPS & 6 & 6 & yes & $-1.671926221$ & $3.75\,10^{-10}$ &
      $0.8773282$ & $0.3559339$ & $(8.46\,10^{-8},0.112)$ \\
    \bottomrule
  \end{tabular}
\end{table}

For TFIM, the structured ansatz is much more accurate than an MPS at the same
nominal $D$, but that comparison gives LETTA the larger transfer bond
$\chi=2D$.  At equal transfer bond, compare VULETTA $D=2$ with VUMPS
$D=4$, or VULETTA $D=3$ with VUMPS $D=6$.  In this small implementation
VUMPS is faster and, at those equal-$\chi$ points, more accurate.

\begin{table}[H]
  \centering
  \caption{Antiferromagnetic Heisenberg snapshot.  Runtime is seconds.}
  \label{tab:heisenberg-benchmark}
  \scriptsize
  \begin{tabular}{@{}llrcrrrrr@{}}
    \toprule
    Method & $D$ & $\chi$ & Conv & $e$ & $|e-e_0|$ &
      $\langle S^z\rangle$ & $\langle S^zS^z\rangle$ & $(\epsilon,\ t)$ \\
    \midrule
    VULETTA & 1 & 2 & yes & $0.212962963$ & $6.56\,10^{-1}$ &
      $1.61\,10^{-7}$ & $0.1666667$ & $(8.54\,10^{-8},0.027)$ \\
    VULETTA & 2 & 4 & no & $-0.407782577$ & $3.54\,10^{-2}$ &
      $1.64\,10^{-11}$ & $-0.1371459$ & $(3.44\,10^{-5},0.425)$ \\
    VULETTA & 3 & 6 & no & $-0.435536167$ & $7.61\,10^{-3}$ &
      $6.01\,10^{-14}$ & $-0.1572787$ & $(4.16\,10^{-5},1.596)$ \\
    \midrule
    VUMPS & 1 & 1 & no & $0.250000000$ & $6.93\,10^{-1}$ &
      $-0.1106$ & $0.01223$ & $(1.414,0.215)$ \\
    VUMPS & 2 & 2 & no & $0.226564773$ & $6.70\,10^{-1}$ &
      $-0.4675$ & $0.23330$ & $(1.414,0.504)$ \\
    VUMPS & 3 & 3 & no & $0.221776800$ & $6.65\,10^{-1}$ &
      $-0.2481$ & $0.04213$ & $(1.399,0.971)$ \\
    VUMPS & 4 & 4 & no & $-62.9335734$ & $6.25\,10^{1}$ &
      $85.0204$ & $-13.8095$ & $(1.414,1.624)$ \\
    VUMPS & 6 & 6 & no & $-0.277443776$ & $1.66\,10^{-1}$ &
      $-0.04963$ & $-0.07270$ & $(1.222,4.088)$ \\
    \bottomrule
  \end{tabular}
\end{table}

\paragraph{Critical warning about the Heisenberg VUMPS rows.}
The present one-site fixed-point update enters an undamped two-cycle and
reaches the iteration limit.  The $D=4$ transfer fixed point is nonpositive,
which produces unphysical observables.  These rows are failure diagnostics,
not variational MPS estimates.  A two-site unit cell and/or a damped mixed
update is needed before using this solver as a Heisenberg baseline.

\subsection{What the tests establish}

The focused tests are deliberately layered:
\begin{enumerate}
  \item exact equality of periodic LETTA and both sparse MPS embeddings;
  \item state preservation, sector isometries, and blockwise center
        identities after conditional canonicalization;
  \item horizontality and strict descent of the conditional direction;
  \item equality of conditional residual/rank with the dense Gram oracle;
  \item gauge invariance of amplitudes, energy, gradient annihilation, and
        tangent residual;
  \item analytic-gradient agreement with real and complex finite
        differences;
  \item rejection of noninjective states and numerically invalid canonical
        retractions;
  \item explicit proof by monkeypatch that the default solver does not build
        the dense tangent Gram matrix;
  \item TFIM convergence, observable references, benchmark serialization,
        VUMPS environments, effective-Hamiltonian Hermiticity and covariance,
        and matrix-free branches.
\end{enumerate}
These tests do not prove global convergence, nor do they make a failed
Heisenberg fixed-point iteration physically valid.  They establish the
algebraic invariants and local derivative/update contracts implemented by the
code.

\section{VULETTA versus VUMPS: exact correspondences and real differences}
\label{sec:comparison}

\begin{table}[htbp]
  \centering
  \caption{Conceptual map between the current implementations.}
  \label{tab:method-map}
  \small
  \begin{tabular}{@{}p{0.20\linewidth}p{0.35\linewidth}p{0.35\linewidth}@{}}
    \toprule
    Concept & VULETTA & VUMPS \\
    \midrule
    Variational tensor
      & $T^{p,q}\in\CC^{D\times D}$
      & $A^s\in\CC^{\chi\times\chi}$ unrestricted \\
    Physical amplitude
      & $\Tr\prod_n T^{s_n,s_{n+1}}$
      & $\Tr\prod_n A^{s_n}$ \\
    Contraction tensor
      & exact sparse embedding $B^s$ of bond $dD$
      & the optimized $A^s$ itself \\
    Gauge
      & $T^{p,q}\mapsto G_p^{-1}T^{p,q}G_q$
      & $A^s\mapsto G^{-1}A^sG$ \\
    Canonical tensors
      & $T_L^{p,q},T_R^{p,q},T_C^{p,q},C_p$
      & $A_L^s,A_R^s,A_C^s,C$ \\
    Isometry
      & one equation for every conditioned sector
      & one left and one right equation \\
    Optimization
      & analytic energy gradient, horizontal projection, metric whitening,
        Armijo retraction
      & center-site and center-bond ground-state eigenproblems, polar matching \\
    Stationarity
      & tangent residual plus canonical residual
      & canonical-consistency plus fixed-point-change residual \\
    \bottomrule
  \end{tabular}
\end{table}

\subsection{The gauge condition really is the MPS condition, but conditioned}

Substituting the exact right-copy embedding into an ordinary MPS gauge shows
why the sector gauge is the correct one.  Define the block-diagonal enlarged
matrix
\begin{equation}
  \mathbb G_{(\alpha,p),(\beta,q)}
  =
  \delta_{pq}(G_p)_{\alpha\beta}.
  \label{eq:enlarged-block-gauge}
\end{equation}
Then
\begin{equation}
  B^s(T')=\mathbb G^{-1}B^s(T)\mathbb G
  \label{eq:embedding-gauge-equivariance}
\end{equation}
holds exactly when $T'^{p,q}=G_p^{-1}T^{p,q}G_q$.  The conditional gauge is
therefore not an analogy added after the fact; it is the restriction of the
ordinary MPS virtual similarity transform to block-diagonal transformations
that preserve the sparse LETTA embedding.

What must \emph{not} be done is to allow a general
$\mathbb G\in\operatorname{GL}(dD,\CC)$.  A generic enlarged transformation
mixes the physical-copy sectors and destroys
$B^s_{(\alpha,p),(\beta,q)}=\delta_{q,s}T^{p,s}_{\alpha\beta}$.  Such a gauge
is legal for the represented MPS but leaves the structured LETTA coordinate
chart.

\subsection{Parameter and tangent counts}

At LETTA bond $D$,
\begin{align}
  N_T&=d^2D^2
  &&\text{complex stored tensor entries},
  \label{eq:letta-storage-count-comparison}\\
  \chi_{\mathrm{transfer}}&=dD,
  \label{eq:letta-transfer-bond}\\
  N_{\mathrm{tan}}&=d(d-1)D^2
  &&\text{generic complex horizontal coordinates}.
  \label{eq:letta-tangent-count}
\end{align}
An ordinary MPS of nominal bond $D$ has
\begin{equation}
  N_A=dD^2,
  \qquad
  N_{\mathrm{tan,MPS}}=(d-1)D^2
  \label{eq:mps-counts-same-D}
\end{equation}
after quotienting its ordinary gauge and scale.  Hence at the same nominal
$D$, LETTA has $d$ times as many stored entries and $d$ times as many generic
tangent coordinates, while its transfer bond is also $d$ times larger.

At the fairer comparison of equal transfer bond $\chi=dD$,
\begin{align}
  N_T&=\chi^2,
  &
  N_A&=d\chi^2,
  \label{eq:equal-transfer-storage}\\
  N_{\mathrm{tan}}^{\mathrm{LETTA}}
    &=\frac{d-1}{d}\chi^2,
  &
  N_{\mathrm{tan}}^{\mathrm{MPS}}
    &=(d-1)\chi^2.
  \label{eq:equal-transfer-tangent}
\end{align}
Thus the structured family uses a factor $1/d$ of the unrestricted MPS
storage and generic tangent dimension at equal contraction bond.  This
compression is also an expressivity restriction: VULETTA cannot in general
reach every MPS of bond $dD$.

\subsection{Why VULETTA does not simply call VUMPS}

VUMPS solves an effective ground-state equation in the full one-site MPS
manifold.  A lowest eigenvector of Eq.~\eqref{eq:vumps-center-eigenproblem}
will generally violate the sparse-copy constraint of
Eq.~\eqref{eq:right-copy-embedding}.  Projecting that eigenvector back to
$T^{p,q}$ would not preserve the variational optimality of the VUMPS step and
would require a separate nonlinear projection.  The present VULETTA solver
instead differentiates the energy \emph{within} the LETTA manifold and
constructs a descent in its quotient tangent space.

\section{Computational cost and current scaling bottlenecks}
\label{sec:complexity}

\subsection{VULETTA}

Let $\chi=dD$.  The current thermodynamic contraction builds a dense transfer
matrix acting on $\chi\times\chi$ matrices:
\begin{equation}
  \mathbf E_A\in\CC^{\chi^2\times\chi^2}.
  \label{eq:dense-transfer-shape}
\end{equation}
Materializing it costs
\begin{equation}
  O(\chi^4)\ \text{memory},
  \label{eq:dense-transfer-memory}
\end{equation}
and a generic dense eigendecomposition or inverse costs
\begin{equation}
  O((\chi^2)^3)=O(\chi^6)
  \ \text{time}.
  \label{eq:dense-transfer-time}
\end{equation}
The analytic gradient loops over all $d^2D^2$ structured parameters and two
real directions for a complex tensor, reusing dense transfer-response
objects.  The conditional tangent construction itself is cheaper: $d$
sector QR decompositions and eigendecompositions of $D\times D$ metric blocks.

Consequently, the default method removes the catastrophic
$O(N_T^2)$ storage of the dense parameter Gram matrix, but it does
\emph{not} yet remove the dense $\chi^2\times\chi^2$ transfer bottleneck.
Exploiting the delta structure in $B^s$ and using matrix-free fixed-point and
reduced-resolvent solvers is the natural next scaling improvement.

\subsection{VUMPS}

VUMPS uses transfer and environment spaces of vector dimension $\chi^2$,
center space of dimension $d\chi^2$, and bond space of dimension $\chi^2$.
Below the thresholds in Table~\ref{tab:vumps-options}, the current code
materializes dense environment/effective matrices.  Above threshold it uses
matrix-free GMRES for environments and a matrix-free Hermitian eigensolver
for the effective Hamiltonians.  This gives VUMPS a materially better
large-$\chi$ route than the present dense VULETTA transfer implementation,
although individual contractions still have the standard MPS polynomial
cost.

\begin{table}[htbp]
  \centering
  \caption{Dominant spaces in the current dense/sparse implementations.}
  \label{tab:complexity-spaces}
  \begin{tabular}{@{}lll@{}}
    \toprule
    Object & VULETTA dimension & VUMPS dimension \\
    \midrule
    transfer vector & $(dD)^2$ & $\chi^2$ \\
    raw variational entries & $d^2D^2$ & $d\chi^2$ \\
    generic quotient tangent & $d(d-1)D^2$ & $(d-1)\chi^2$ \\
    center-site eigenproblem & not used & $d\chi^2$ \\
    center-bond eigenproblem & not used & $\chi^2$ \\
    \bottomrule
  \end{tabular}
\end{table}

\section{The VUMPS implementation used for comparison}
\label{sec:vumps}

\subsection{Ordinary one-site uniform MPS and its gauge}

The neighboring module \path{pyqed/_vumps} optimizes an unrestricted one-site
uniform MPS tensor
\begin{equation}
  A=\{A^s\}_{s=0}^{d-1},
  \qquad A^s\in\CC^{\chi\times\chi},
  \label{eq:ordinary-mps-tensor}
\end{equation}
with periodic amplitude
\begin{equation}
  \Psi_A(s_1,\ldots,s_N)
  =
  \Tr\!\left(A^{s_1}A^{s_2}\cdots A^{s_N}\right).
  \label{eq:ordinary-mps-amplitude}
\end{equation}
Its familiar gauge freedom is
\begin{equation}
  A^s\longmapsto G^{-1}A^sG,
  \qquad G\in\operatorname{GL}(\chi,\CC).
  \label{eq:ordinary-mps-gauge}
\end{equation}
Equation~\eqref{eq:letta-gauge} is precisely the conditional version of this
MPS gauge: the single $G$ is replaced by a family $\{G_p\}$ attached to the
physical label carried by the enlarged virtual leg.

The MPS transfer actions in \code{state.py} are
\begin{align}
  \mathcal T_R^A(X)&=\sum_s A^sX A^{s\dagger},
  \label{eq:vumps-right-transfer}\\
  \mathcal T_L^A(X)&=\sum_s A^{s\dagger}X A^s.
  \label{eq:vumps-left-transfer}
\end{align}
These are the same two maps as $\E_A$ and $\E_A^\dagger$ in
Eq.~\eqref{eq:transfer-maps}.

\subsection{Mixed-canonical MPS form}

The \code{CanonicalMPS} container stores $A_L$, $A_R$, a center matrix $C$,
and optionally an independently optimized center tensor $A_C$.  At a fixed
point they obey
\begin{equation}
  A_C^s=A_L^sC=CA_R^s
  \qquad\text{for every }s,
  \label{eq:vumps-center-relation}
\end{equation}
together with
\begin{align}
  \sum_s A_L^{s\dagger}A_L^s&=\Id,
  \label{eq:vumps-left-isometry}\\
  \sum_s A_R^sA_R^{s\dagger}&=\Id,
  \label{eq:vumps-right-isometry}\\
  \norm{C}_F&=1.
  \label{eq:vumps-center-normalization}
\end{align}
The two bond density matrices are
\begin{equation}
  \rho_L=\frac{CC^\dagger}{\Tr(CC^\dagger)},
  \qquad
  \rho_R=\frac{C^\dagger C}{\Tr(C^\dagger C)}.
  \label{eq:vumps-bond-densities}
\end{equation}

The routine \code{canonicalize} begins with dominant left and right fixed
points $l,r$ of a raw $A$.  After dividing $A$ by the square root of its
transfer radius, it computes positive roots and forms
\begin{align}
  A_L^s&=l^{1/2}A_0^s l^{-1/2},
  \label{eq:vumps-canonical-left}\\
  A_R^s&=r^{-1/2}A_0^s r^{1/2},
  \label{eq:vumps-canonical-right}\\
  C&=\frac{l^{1/2}r^{1/2}}
          {\norm{l^{1/2}r^{1/2}}_F}.
  \label{eq:vumps-canonical-center}
\end{align}
These are the unconditioned analogues of
Eqs.~\eqref{eq:TL}--\eqref{eq:conditional-centers}.
Both implementations reject substantially negative or rank-deficient fixed
points rather than silently pseudoinverting a noninjective state.

\subsection{Observables in the VUMPS representation}

The code treats $A_L$ as the definite physical uMPS representative.  If
$\rho$ is the normalized right fixed point of $\mathcal T_R^{A_L}$, then
\begin{equation}
  \langle o\rangle
  =
  \sum_{s,u}o_{su}
  \Tr\!\left(
    \rho\,A_L^{s\dagger}A_L^u
  \right)
  \label{eq:vumps-one-site-observable}
\end{equation}
and
\begin{equation}
  e
  =
  \sum_{s,t,u,v}h_{stuv}
  \Tr\!\left[
    \rho\,
    A_L^{t\dagger}A_L^{s\dagger}A_L^uA_L^v
  \right].
  \label{eq:vumps-energy}
\end{equation}
During an unfinished VUMPS iteration the stored $A_C$ can be a provisional
eigensolver output not exactly equal to $A_LC$ or $CA_R$.  The energy and
observables deliberately use $A_L$, not that provisional center.

\subsection{Semi-infinite Hamiltonian environments}

Define the uncentered left and right two-site source matrices
\begin{align}
  S_L
  &=
  \sum_{s,t,u,v}h_{stuv}
  A_L^{t\dagger}A_L^{s\dagger}A_L^uA_L^v,
  \label{eq:vumps-left-source}\\
  S_R
  &=
  \sum_{s,t,u,v}h_{stuv}
  A_R^uA_R^vA_R^{t\dagger}A_R^{s\dagger}.
  \label{eq:vumps-right-source}
\end{align}
The source builders Hermitize these matrices.  The energy density is
$e=\Tr(\rho_L S_L)$ on the definite $A_L$ state, and the environment builder
first subtracts $e\Id$.

For either orientation, let $\mathcal T$ be the relevant canonical transfer
map and let
\begin{equation}
  \Phi_\rho(X)=\Tr(\rho X).
  \label{eq:rho-functional}
\end{equation}
The singular identity component of $I-\mathcal T$ is regularized by solving
\begin{equation}
  \left[
    I-\mathcal T+|\Id)\Phi_\rho
  \right]H
  =
  S-\Id\,\Phi_\rho(S).
  \label{eq:vumps-environment-equation}
\end{equation}
After solving, the implementation enforces
\begin{equation}
  H\leftarrow\tfrac12(H+H^\dagger),
  \qquad
  H\leftarrow H-\Id\,\Phi_\rho(H).
  \label{eq:vumps-environment-gauge}
\end{equation}
For vectorized dimension $\chi^2\leq256$ it explicitly builds and solves the
dense linear system.  Above that threshold it applies the same operator
matrix-free with GMRES.  The two solutions are denoted $H_L$ and $H_R$.

\subsection{Effective center-site Hamiltonian}

For a trial center tensor $X=\{X^x\}$, the implemented linear action is
\begin{align}
  (\mathcal H_{A_C}X)^x
  ={}&
  H_LX^x+X^xH_R
  \notag\\
  &+\sum_{s,u,v}
    h_{s x u v}\,
    A_L^{s\dagger}A_L^uX^v
  \notag\\
  &+\sum_{t,u,v}
    h_{x t u v}\,
    X^uA_R^vA_R^{t\dagger}.
  \label{eq:vumps-effective-center}
\end{align}
The first local term places $X$ on the right site of the interaction; the
second places it on the left site.  The new center is the normalized lowest
eigenvector,
\begin{equation}
  A_C^{\mathrm{new}}
  =
  \argmin_{\norm{X}_F=1}
  \langle X,\mathcal H_{A_C}X\rangle_F.
  \label{eq:vumps-center-eigenproblem}
\end{equation}

\subsection{Effective center-bond Hamiltonian}

For a trial bond matrix $Y$, the action is
\begin{align}
  \mathcal H_C(Y)
  ={}&
  H_LY+YH_R
  \notag\\
  &+\sum_{s,t,u,v}h_{stuv}
  A_L^{s\dagger}A_L^u
  Y
  A_R^vA_R^{t\dagger},
  \label{eq:vumps-effective-bond}
\end{align}
and
\begin{equation}
  C^{\mathrm{new}}
  =
  \argmin_{\norm{Y}_F=1}
  \langle Y,\mathcal H_C(Y)\rangle_F.
  \label{eq:vumps-bond-eigenproblem}
\end{equation}
For a vector space of dimension at most
\code{dense_eigensolver_threshold} $=256$, the code materializes the
Hermitian matrix and calls a dense eigensolver.  Otherwise it uses a
matrix-free smallest-algebraic-eigenvalue solve.  Each eigenvector is phase
aligned with the previous center before normalization.

\subsection{Polar gauge matching}

The two eigenproblems optimize $A_C$ and $C$ independently, so the solver must
recover isometric $A_L$ and $A_R$.  Let
\begin{align}
  A_{C,L}&=\reshape(A_C,(d\chi)\times\chi),
  \label{eq:left-center-reshape}\\
  A_{C,R}&=\reshape(A_C,\chi\times(d\chi)),
  \label{eq:right-center-reshape}
\end{align}
and let
\begin{equation}
  \operatorname{polar}(M)=UV^\dagger
  \quad\text{when}\quad M=U\Sigma V^\dagger .
  \label{eq:polar-map}
\end{equation}
The implemented gauge match is
\begin{align}
  A_L&=
  \reshape\!\left[
    \operatorname{polar}(A_{C,L})
    \operatorname{polar}(C)^\dagger
  \right],
  \label{eq:vumps-polar-left}\\
  A_R&=
  \reshape\!\left[
    \operatorname{polar}(C)^\dagger
    \operatorname{polar}(A_{C,R})
  \right].
  \label{eq:vumps-polar-right}
\end{align}
This is not an energy line search.  It is a fixed-point map built from the
two effective eigenvectors.

\subsection{VUMPS residuals and iteration}

After gauge matching, define
\begin{align}
  \epsilon_{\mathrm{can}}^{\mathrm{V}}
  &=
  \max\left\{
    \norm{A_C-A_LC}_F,\,
    \norm{A_C-CA_R}_F
  \right\},
  \label{eq:vumps-canonical-residual}\\
  \epsilon_{\mathrm{fp}}^{\mathrm{V}}
  &=
  \max\left\{
    \norm{A_C^{\mathrm{new}}-A_C^{\mathrm{old}}}_F,\,
    \norm{C^{\mathrm{new}}-C^{\mathrm{old}}}_F
  \right\},
  \label{eq:vumps-fixedpoint-residual}\\
  \epsilon_{\mathrm{VUMPS}}
  &=
  \max\left\{
    \epsilon_{\mathrm{can}}^{\mathrm{V}},
    \epsilon_{\mathrm{fp}}^{\mathrm{V}}
  \right\}.
  \label{eq:vumps-total-residual}
\end{align}
The current one-site VUMPS loop repeats environment construction, the two
lowest-eigenpair solves, polar matching, and residual evaluation until
$\epsilon_{\mathrm{VUMPS}}$ reaches the requested tolerance.

\begin{table}[htbp]
  \centering
  \caption{Defaults in \code{VUMPSOptions}.}
  \label{tab:vumps-options}
  \begin{tabular}{@{}lll@{}}
    \toprule
    Option & Default & Meaning \\
    \midrule
    \code{max_iterations} & $100$ & fixed-point iterations \\
    \code{tolerance} & $10^{-8}$ & total residual target \\
    \code{eigensolver_tolerance} & $10^{-10}$ & iterative eigenproblem target \\
    \code{dense_eigensolver_threshold} & $256$ & vector-space cutoff \\
    \code{environment_tolerance} & $10^{-12}$ & GMRES target \\
    \code{dense_environment_threshold} & $256$ & vectorized bond-space cutoff \\
    \code{verbosity} & $0$ & progress printing \\
    \bottomrule
  \end{tabular}
\end{table}

\section{Final synthesis}

The two solvers share the same transfer-fixed-point and mixed-canonical
language, but operate on different manifolds.  VUMPS uses the full MPS gauge
and unrestricted center eigenvectors.  VULETTA preserves the nearest-neighbor
leg tying by restricting that MPS gauge to the block-diagonal family
$\{G_p\}$, and it optimizes only the resulting horizontal structured tangent.
This is the precise sense in which the current implementation now uses the
MPS-like gauge condition fully.

\end{document}
